\documentclass[../../../include/open-logic-section]{subfiles}

\begin{document}

\olfileid{sth}{card-arithmetic}{fix}
\olsection{$\aleph$-स्थिरबिंदू}

\olref[spine][]{chap} मध्ये आपण असे सुचवले होते की प्रतिस्थापनाला समर्थनाची गरज आहे, कारण त्याच्यामुळे श्रेणीची उंची बरीच मोठी असावी लागते. संचांकांचे काही अंकगणित केल्यानंतर या उंचीचे एक छोटे उदाहरण देता येईल.

स्पष्टच $0 < \aleph_0$, $1 < \aleph_1$, $2 < \aleph_2$\ldots आणि प्रत्येक पावलागणिक आकारमानातील फरक आणखी \emph{वाढतो}. म्हणून प्रत्येक क्रमसूचक संख्या $\kappa$ साठी $\kappa< \aleph_\kappa$ असा तर्क करण्याचा मोह होतो.

पण $\ZFC$ स्वीकारल्यास हा तर्क \emph{खोटा} ठरतो. खरे तर \emph{$\aleph$-स्थिरबिंदू} अस्तित्वात असल्याचे सिद्ध करता येते; म्हणजे $\kappa=\aleph_\kappa$ असलेले संचांक $\kappa$ आहेत.

\begin{prop}\ollabel{alephfixed}
एक $\aleph$-स्थिरबिंदू अस्तित्वात आहे.
\end{prop}

\begin{proof}
पुनरावर्तनाने पुढील व्याख्या करा:
\begin{align*}
	\kappa_0 &= 0\\
	\kappa_{n+1} &= \aleph_{\kappa_n}\\
	\kappa&= \bigcup_{n < \omega}\kappa_n
\end{align*}
\olref[cardinals][classing]{unioncardinalscardinal} नुसार $\kappa$ हा संचांक आहे. अनुक्रम वाढत असून त्याची मूल्ये त्यांच्या युतीकरणात सहअंतिम असल्यामुळे:
\[
	\kappa= \bigcup_{n < \omega} \kappa_{n+1} =
	\bigcup_{n < \omega}\aleph_{\kappa_n} =
	\bigcup_{\alpha < \kappa}\aleph_\alpha = \aleph_\kappa
\]
%	रचनेनुसार $\kappa$ हा प्रत्येक
%	$\kappa_n$ पेक्षा मोठा असलेला लघुतम संचांक आहे. त्यामुळे $\aleph_\kappa$ हा प्रत्येक
%	$\aleph_{\kappa_n}$ पेक्षा मोठा असलेला लघुतम संचांक, आणि म्हणून प्रत्येक
%	$\kappa_{n+1}$ पेक्षा मोठा असतो. पण तितक्याच रीतीने $\kappa$ हा प्रत्येक
%	$\kappa_{n+1} = \aleph_{\kappa_n}$ पेक्षा मोठा असलेला लघुतम संचांक आहे. म्हणून $\kappa=
%	\aleph_\kappa$.
\end{proof}

आत्ताच रचलेल्या नेमक्या $\aleph$-स्थिरबिंदूविषयी बूलोस यांनी एक लेख लिहिला होता. लेखाच्या सुरुवातीला $\kappa$ चे अस्तित्व नोंदवल्यावर ते म्हणाले:
\begin{quote}
	संच उपपत्तीशी यापूर्वी परिचय नसलेल्यांच्या दृष्टीने [$\kappa$ ही] \emph{बरीच मोठी} संख्या आहे. ती इतकी मोठी आहे की, निदान $\kappa$ वस्तू आहेत असे ज्यातील एखादे विधान सांगते, अशा कोणत्याही उपपत्तीच्या सत्यतेविषयी प्रश्न उपस्थित करावा असे मला वाटते.
	\citep[p.~257]{Boolos2000}
\end{quote}
आणि शेवटी त्यांनी आपला लेख पुढील प्रश्नाने संपवला:
\begin{quote}
	कथेच्या सुरुवातीला काहीही असले, तरी इथवर येईपर्यंत चाके नुसतीच फिरत आहेत आणि प्रत्यक्षात जे काही आहे त्याचे वर्णन आपण आता ऐकतच नाही, अशी शंका आपल्याला येते का?
	\citep[p.~268]{Boolos2000}
\end{quote}
प्रत्यक्षात “जे काही आहे” त्याच्या पलीकडे आपण गेलो असू, तर दोषाचे बोट थेट प्रतिस्थापनाकडे दाखवावे लागेल. कारण आपला पुरावा याच स्वयंसिद्धकामुळे चालतो. असे असल्यास, काही दशकांपूर्वी केलेल्या प्रतिस्थापनाचे “कोणतेही अनिष्ट” परिणाम नाहीत या दाव्याचा बूलोस यांना पुनर्विचार करावा लागेल (पाहा \olref[replacement][extrinsic]{sec}).

पण $\kappa$ चे अस्तित्व इतके वाईट आहे का? येथे रसेलचा \emph{ट्रिस्ट्रम शॅंडी विरोधाभास} पाहणे उपयुक्त ठरेल. ट्रिस्ट्रम शॅंडी आपल्या जीवनाची नोंद दैनंदिनीत करतो, पण एका दिवसाची नोंद करायला त्याला एक वर्ष लागते. प्रत्येक वर्षागणिक तो अधिकाधिक मागे पडतो: एका वर्षानंतर त्याने केवळ एका दिवसाची नोंद केली आहे आणि नोंद न झालेले 364 दिवस जगला आहे; दोन वर्षांनंतर केवळ दोन दिवसांची नोंद केली आहे आणि नोंद न झालेले 728 दिवस जगला आहे; तीन वर्षांनंतर केवळ तीन दिवसांची नोंद केली आहे आणि नोंद न झालेले 1092 दिवस जगला आहे \dots\footnote{येथे लीप वर्षांचा विचार केलेला नाही.} तरी ट्रिस्ट्रम \emph{अमर} असेल, तर तो प्रत्येक दिवसाची नोंद करू शकेल, कारण आपल्या आयुष्याच्या $n$ व्या वर्षी तो $n$ व्या दिवसाची नोंद करेल. म्हणून “काळाच्या अखेरीस” त्याची दैनंदिनी पूर्ण असेल.

आता, $\kappa$ येईपर्यंत $\alpha$ हा $\aleph_\alpha$ पेक्षा लहान असतो---आणि वाढत्या प्रमाणात, अगदी हतबल करणाऱ्या प्रमाणात लहान असतो---पण त्या बिंदूवर आपण गाठतो आणि $\kappa
= \aleph_\kappa$ मिळते, या विचारापेक्षा वरील गोष्ट इतकी वेगळी का आहे?

हा प्रश्न बाजूला ठेवून, आणि $\ZFC$ स्वीकारतो असे गृहीत धरून, स्थिरबिंदू रचनांविषयीच्या आणखी काही रंजक गोष्टींनी हे प्रकरण संपवू. पुढील तीन निष्कर्षांचा सहज अर्थ असा की श्रेणीची उंची आणि रुंदी समान असलेला एक क्षुल्लक नसलेला टप्पा आहे:

\begin{prop}\ollabel{bethfixed}
एक $\beth$-स्थिरबिंदू अस्तित्वात आहे; म्हणजे $\kappa=
\beth_\kappa$ असलेला $\kappa$ आहे.
\end{prop}

\begin{proof}
\olref{alephfixed} प्रमाणे, “$\aleph$” ऐवजी “$\beth$” वापरा.
\end{proof}

\begin{prop}\ollabel{stagesize}
$\card{V_{\omega+\alpha}} = \beth_{\alpha}$. जर $\omega \ordtimes
\omega \leq \alpha$ असेल, तर $\card{V_\alpha} = \beth_\alpha$.
\end{prop}

\begin{proof}
पहिला दावा साध्या अनंतपार विगमनाने मिळतो. दुसरा दावा यावरून मिळतो की $\omega \ordtimes \omega \leq \alpha$ असेल, तर $\omega + \alpha = \alpha$. हे सिद्ध करण्यासाठी \olref[ord-arithmetic][]{chap} मधील क्रमसूचक अंकगणिताची तथ्ये वापरू. आधी पाहा की $\omega \ordtimes \omega = \omega \ordtimes (1 \ordplus \omega) =
(\omega \ordtimes 1) \ordplus (\omega\ordtimes\omega) = \omega
\ordplus (\omega \ordtimes \omega)$. आता $\omega \ordtimes \omega
\leq \alpha$ असल्यास, म्हणजे एखाद्या $\beta$ साठी $\alpha = (\omega\ordtimes\omega) \ordplus \beta$ असल्यास, $\omega \ordplus \alpha = \omega \ordplus
((\omega \ordtimes \omega) \ordplus \beta) = (\omega \ordplus (\omega
\ordtimes \omega)) \ordplus \beta = (\omega \ordtimes \omega) \ordplus
\beta = \alpha$.
\end{proof}

\begin{cor}
$\card{V_\kappa} = \kappa$ असलेला $\kappa$ अस्तित्वात आहे.
\end{cor}

\begin{proof}
\olref{bethfixed} नुसार $\kappa$ हा एक $\beth$-स्थिरबिंदू असू द्या. स्पष्टच $\omega \ordtimes \omega < \kappa$. म्हणून \olref{stagesize} नुसार $\card{V_\kappa} =
\beth_\kappa= \kappa$.
\end{proof}

$V_\kappa$ च्या खाली जितके टप्पे आहेत, तितकेच !!{element}s $V_\kappa$ मध्ये आहेत. त्यामुळे सहज अर्थाने $V_\kappa$ ची रुंदी त्याच्या उंचीइतकी आहे. हे ट्रिस्ट्रम शॅंडीच्या गोष्टीसारखेच आहे: \emph{घातसंच} घेऊन आपण एका टप्प्यापासून पुढच्या टप्प्यावर जातो, आणि प्रत्येक पावलावर श्रेणी \emph{खूप} मोठी होते. पण “अखेरीस”, म्हणजे $\kappa$ टप्प्यावर, श्रेणीची रुंदी तिच्या उंचीला गाठते.

प्रश्न पडू शकतो: \emph{श्रेणीची रुंदी तिच्या उंचीशी किती वेळा जुळते?} उत्तर असे: \emph{जितक्या क्रमसूचक संख्या आहेत, तितक्या वेळा.} पण याचे थोडे स्पष्टीकरण हवे.

पुढीलप्रमाणे $\tau$ हे पद परिभाषित करू. कोणत्याही $A$ साठी:
\begin{align*}
	\tau_0(A) & \defis \cardsucc{\card{A}}\\
	\tau_{n+1}(A) & \defis \beth_{\tau_n(A)}\\
	\tau(A) & \defis \bigcup_{n < \omega}\tau_n(A)
\intertext{\olref{bethfixed} प्रमाणे, कोणत्याही $A$ साठी $\tau(A)$ हा
$\beth$-स्थिरबिंदू आहे. सुरुवात उत्तरवर्ती संचांकापासून केल्यामुळे $\card{A} < \tau(A)$ मिळते. म्हणून पुढील पुनरावर्ती व्याख्या पाहा:}
	W_0 &\defis 0\\
	W_{\alpha + 1} & \defis \tau(W_\alpha)\\
	W_\alpha & \defis \bigcup_{\beta < \alpha} W_\beta
	\text{, जेव्हा $\alpha$ सीमा क्रमसूचक संख्या असते}
\end{align*}
ही रचना सर्व क्रमसूचक संख्यांसाठी परिभाषित आहे. सहज अर्थाने $W$ हे क्रमसूचक संख्यांपासून $\beth$-स्थिरबिंदूंकडे जाणारे “!!a{injection}” आहे; यातील शून्यव्या मूल्याचा अपवाद स्पष्ट ठेवावा. आणि आधीप्रमाणेच प्रत्येक $\alpha$ साठी $V_{W_\alpha}$ ची रुंदी त्याच्या उंचीइतकी आहे. शून्यव्या टप्प्यावर दोन्ही आकार शून्य आहेत; पुढील टप्प्यांवर हे बेथ-स्थिरबिंदूंच्या गुणधर्मामुळे मिळते.

\end{document}
